\chapter*{Introducci\'on}
% Aquí hablaremos de la convención actual de la velocidad de la luz como constante, de algunas implicaciones de esto como el uso del efecto doppler para medir distancias, se dejará entre abierta la idea de que el paradigma puede cambiar ...
% Se especificará los temas a tratar en cada capítulo


% ### Parte 1: Contextualización y Paradigma Actual
% *   **Párrafo 1 (El postulado de la velocidad de la luz constante):**
%     *   *De qué hablar:* Introduce el postulado fundamental de la Teoría de la Relatividad Especial de Albert Einstein: la constancia de la velocidad de la luz (\\(c\\)) en el vacío para cualquier observador inercial.
%     *   *Qué mencionar:* Las enormes implicaciones que este supuesto tiene en la física contemporánea y en la metrología (por ejemplo, la definición actual del metro o el uso clásico del efecto Doppler para calcular distancias astronómicas).

En 1905, Albert Einstein publicó la Teoría de la Relatividad Especial, la cual revolucionó la física con dos postulados simples: las leyes de la física son iguales para todos los observadores incerciales, y la luz viaja a una velocidad constante en el vacío. Las implicaciones de estos fueron enormes y llegaron a muchas ramas de la ciencia; por ejemplo, permitió calcular las velocidades de objetos astronómicos mediante el efecto Doppler. Este efecto ya era conocido en ese entonces por su versión sonora, donde si la fuente se acerca o se aleja se puede percibir un cambio en la frecuencia del sonido, haciéndolo más agudo o más grave. Este fenómeno fue fácilmente extrapolado a la luz, donde ahora el cambio de la frecuencia se ve como un cambio en el color: un corrimiento hacia el rojo o al azul. Al conocer la velocidad de la luz, la frecuencia de la señal emitida y la diferencia de dicha frecuencia con respecto a la detectada, es posible conocer la velocidad de desplazamiento de la fuente en dirección radial. Otra implicación de la velocidad de la luz constante está en la actual definición del metro, pues se estableció como la distancia que recorre la luz en un cierto tiempo, convirtiéndolo así en una unidad de distancia que se puede usar universalmente.

% *   **Párrafo 2 (La noción del cambio de paradigma):**
%     *   *De qué hablar:* Plantea de manera elegante e intelectual que, históricamente, los dogmas de la física han cambiado ante nueva evidencia experimental.
%     *   *Qué mencionar:* Introduce la idea de que el vacío podría no ser un escenario pasivo, sino un "medio energético cósmico" con densidades de energía (gravitacionales y electromagnéticas) variables de un punto a otro del universo, lo cual sugiere la posibilidad de que la velocidad de la luz fluctúe según las propiedades de su entorno local.

La física ha tomado la relatividad de Einstein como una base sobre la cual construir y dar explicación a nuevos fenomenos por más de un siglo, y esto, aunque cómodo, no es lo que ocurre historicamente en la ciencia. El método cientifico nos lleva constantemente a repasar y reescribir las explicaciones basicas cuando llega un nuevo fenomeno que ya no se puede describir con esas reglas, esos cambios de paradigma son necesarios para el avance del conocimiento humano. Con ese espiritu, este trabajo busca indagar en la idea de un espacio cósmico que no está del todo vacio, sino que funciona como un medio energético capaz de hacer que la velocidad de la luz fluntue según las propiedades de su entorno local, y con ello dar una explicación a algunos de los más fascinantes fenómenos de la astronomía moderna.
 
% ### Parte 2: Antecedentes del Trabajo (Soporte Físico y Teórico)
% *   **Párrafo 3 (La Hipótesis de Céspedes-Curé - HCC):**
%     *   *De qué hablar:* Presenta la hipótesis central de tu investigación propuesta por el Prof. Jorge Céspedes-Curé en su libro *Einstein on Trial...*.
%     *   *Qué mencionar:* Define matemáticamente la relación fundamental \\(c = k\sqrt{\rho}\\), donde la velocidad de la luz depende directamente de la densidad de energía total del espacio (\\(\rho\\)) en un punto determinado.
% *   **Párrafo 4 (Soporte Astrofísico I: La Anomalía de las Sondas Pioneer):**

La hipótesis medular del presente trabajo es una propuesta por Jorge Céspede-Curé en 2002, en su libro \textit{Einstein on Trial or Metaphysical Principles of Natural Philosophy} \cite{bib:el-libro}, que dice que la velocidad de la luz $c$ depende la de densidad de energía total del espacio $\rho$ como  $c = k / \sqrt{\rho}$, aquí la $\rho$ es la suma de las contribuciones de los campos electromagnético y gravitacional, y $k$ una constante. Este nuevo enfoque ya permitió explicar en 2007 la anomalía de las sondas Pioneer \cite{Cisnotconstant} y en 2020 la anomalía en la maniobra de asistencia gravitacional para los casos de sobrevuelo cercano \cite{Greavesflyby}. La primera, Pioneer Anomaly, consiste en una desaceleración de origen desconocido que presentaron las sondas Pioneer 10 y 11 al salir del sistema solar; la explicación propuesta es que la disminución del tiempo de viaje de la señal, debido a una densidad de enería gravitacional progresivamente menor a medida que la nave se aleja, provoca una desaceleración aparente cuando se uiliza el efecto doppler como medidor de la velocidad. El segundo caso, conocido como Flyby anomaly, consiste en desfases inesperados en las energías cinéticas de las naves entre la entrada y la salida de la maniobra, esto se llegó a observar en varias sondas como en Cassini y Galileo; la explicación propuesta pasa de nuevo por el efecto doppler pero tambien por la asimetría entre los puntos de medida (entrada y salida) y el sistema Tierra-Sol, que origina valores de $\rho$ diferentes en casa punto. 



%     *   *De qué hablar:* Expón el primer gran antecedente de soporte astrofísico analizado por el Dr. Eduardo Greaves (2007).
%     *   *Qué mencionar:* La "Anomalía de las Pioneer" observada por la NASA (una desaceleración no explicada en las sondas Pioneer 10 y 11 al salir del sistema solar) y cómo la HCC propone que esta desaceleración es aparente, causada por el aumento del tiempo de viaje de la señal debido a una disminución de la densidad de energía gravitacional solar a medida que la nave se aleja.
% *   **Párrafo 5 (Soporte Astrofísico II: La Anomalía de Sobrevuelo o "Flyby Anomaly"):**
%     *   *De qué hablar:* Expón el segundo gran antecedente astrofísico publicado por Greaves y colaboradores en 2020.
%     *   *Qué mencionar:* La anomalía orbital registrada en las maniobras de asistencia gravitacional de sondas de la NASA al acercarse a la Tierra (desfases inesperados en la frecuencia Doppler entre la entrada y la salida de la nave) y cómo la asimetría de la densidad de energía gravitacional local del sistema Tierra-Sol resuelve matemáticamente el misterio sin recurrir a "nueva física" gravitacional.
% *   **Párrafo 6 (Soporte Terrestre: Análisis de medidas históricas de \\(c\\)):**
%     *   *De qué hablar:* Describe cómo la HCC también arroja luz sobre las sutiles variaciones estacionales o locales reportadas históricamente en los experimentos terrestres de medición de la velocidad de la luz.


 %En esa direccion ya Rambdani \cite{ramdani} hizo notar una aparente relación entre la velocidad de la luz medida de forma experimental y la gravedad local en el lugar de la medición.



% ### Parte 3: Planteamiento del Problema (La Necesidad de Precisión)
Siendo la \ac{HCC} una idea poco explorada deforma directa,  %le queda aun mucho por demostrar, 
el presente trabajo busca hacer un aporte en su verificación mediante el diseño, caracterización y evaluación de metodologías experimentales, que sean capaces de registrar variaciones en la \ac{c} inducidas por la densidad de energía gravitacional en la superficie terrestre. De acuerdo con las predicciones teóricas de la \ac{HCC}, las diferencias esperadas en el valor de c entre localidades con potenciales gravitacionales contrastantes (como el campus de Sartenejas de la \ac{USB}, el Teleférico de Caracas y Carmen de Uria) se manifiestan únicamente a partir de la novena cifra significativa. Esto impone un desafío metrológico monumental: para corroborar o refutar empíricamente la hipótesis, cualquier sistema de medición debe poseer una precisión límite inferior a $\pm 1$ m/s. 

Para abordar este reto, en esta investigación se plantearon y desarrollaron dos intentos experimentales fundamentados en principios físicos independientes:
El primer intento (Experimento I: medida directa): Consistió en la caracterización directa de la velocidad de propagación de un haz electromagnético en el espacio libre. Se implementó un montaje basado en la modulación en alta frecuencia de un láser de He-Ne y la medición del desfase temporal ($\Delta t$) de la señal sinusoidal en función de la diferencia de camino óptico ($\Delta x$) utilizando un osciloscopio digital. El planteamiento del problema en esta vía radica en determinar si, mediante el refinamiento del alineamiento óptico y la optimización de algoritmos de regresión en Python, es factible alcanzar con instrumentación convencional las resoluciones de escala de decenas de nanómetros y de femtosegundos que se requieren para resolver la novena cifra de la velocidad de la luz.

El segundo intento (Experimento II: medida indirecta): Surgió como una alternativa electrodinámica para inferir la variación de la velocidad de la luz a través de las propiedades dieléctricas del vacío ($\epsilon_0$) mediante la relación de Maxwell $c=1/ \sqrt{\mu_0\epsilon_0}$. Para ello, se diseñó y construyó un equipo físico basado en circuitos resonantes RLC (evaluando configuraciones en serie y paralelo) acoplados a módulos de generación estable y estabilización térmica activa. La problemática en este segundo enfoque consiste en controlar el factor de calidad Q del resonador y mitigar el ruido térmico y las capacidades parásitas del circuito, con el fin de verificar si es posible estabilizar y caracterizar el pico de la frecuencia de resonancia ($f_0$) con la sensibilidad infinitesimal necesaria para registrar la influencia del entorno gravitacional.


% ### Parte 4: Justificación, Importancia e Hipótesis Implícita

A partir de la complementariedad de ambos esquemas de medición, se postula la hipótesis implícita del presente proyecto de grado: si la densidad de energía gravitacional local altera las propiedades del vacío de acuerdo con la relación teórica de Céspedes-Curé, entonces la velocidad de propagación de un haz de luz modulado medida por vía directa óptica, así como la frecuencia de resonancia de un circuito tanque RLC blindado y térmicamente estable, exhibirán alteraciones sistemáticas y medibles de manera correlacionada con el gradiente de potencial gravitatorio terrestre al ser evaluadas en diferentes altitudes de la superficie de la Tierra.

% *   **Párrafo 11 (La Hipótesis Implícita):**
%     *   *De qué hablar:* Formula de forma implícita (como lo pide la norma USB) la hipótesis que guía el proyecto.
%     *   *Qué colocar (redactar algo así):* *"Si las variaciones en la densidad de energía gravitacional local modifican la permitividad eléctrica del espacio, entonces un sistema de circuito resonante RLC operando bajo riguroso control ambiental exhibirá un corrimiento medible en su frecuencia de resonancia al ser trasladado entre locaciones geográficas con diferentes potenciales gravitacionales, reflejando de manera indirecta un cambio en la velocidad de la luz en la Tierra"*.

% ---

% ### Parte 5: Objetivos del Trabajo (Sección Estructurada)
% *(Debe presentarse como un apartado formal e independiente dentro de la introducción, tal como exige la DEP)*

% *   **Objetivo General:**
%     *   Desarrollar un sistema de medición experimental indirecto y de alta precisión, basado en circuitos resonantes RLC, para verificar los efectos de la densidad de energía gravitacional sobre la velocidad de la luz en la superficie terrestre, evaluando cuantitativamente la validez de la Hipótesis de Céspedes-Curé (HCC).
% *   **Objetivos Específicos:**
%     1.  Diseñar y simular un circuito resonante RLC (evaluando configuraciones serie y paralelo) optimizado para maximizar el factor de calidad \\(Q\\), obteniendo una curva de resonancia lo suficientemente estrecha para detectar corrimientos de frecuencia mínimos.
%     2.  Desarrollar una arquitectura electrónica modular compuesta por un módulo generador estable (oscilador AD9833 controlado por Arduino), un módulo de lectura y un módulo de estabilización térmica activa para mitigar las variables ambientales.
%     3.  Construir físicamente los prototipos experimentales en placas de circuito impreso (PCB) para reducir efectos parásitos y ruidos que comprometan las mediciones.
%     4.  Implementar un software de procesamiento de datos en Python que automatice el barrido de frecuencias y ejecute un ajuste matemático robusto mediante Regresión Ortogonal de Distancias (ODR) para calcular con precisión infinitesimal la frecuencia de resonancia y sus incertidumbres.
%     5.  Evaluar la precisión lograda por el sistema, proponer un protocolo de mediciones geográficas y plantear alternativas tecnológicas avanzadas (como la plataforma Red Pitaya o acoplamiento por inducción) para futuras fases experimentales.

Para cumplir lo anterior este trabajo se organiza con cuatro capítulos principales. El capítulo I es una explicación detallada de la \ac{HCC} y su origen, las anomalías de las Pioneer y Flyby, y desarrolla el cálculo de las predicciones de c en cada localidad. El capítulo II expone la implementación del experimento I de medida directa y los resultados obtenidos. En el capítulo III se explica la bases teóricas en el marco de la electrodinámica, la arquitectura y construccion de dos versiones del experimento II, junto con los resultados y propuestas de mejoras. El capítulo IV plantea dos experimentos adicionales. Se incluyen ademas 4 apendices con mayor detalle sobre el experimento II en cuanto a la parte de programación informatica (apéndice A), el tratamiento estadístico de los datos (apéndice B), montaje electrónico (apéndice C) y matemática involucrada (apéndice D).

% ### Parte 6: Estructura del Libro (Transición Final)
% *   **Párrafo 12 (Organización capitular del informe de tesis):**
%     *   *De qué hablar:* Describe brevemente al lector la estructura que tendrá tu libro de grado para facilitar su lectura (esta práctica la puedes observar idéntica en los informes aprobados de Juan Guerrero y Celeste Méndez).
%     *   *Ejemplo de flujo:* Señala que el **Capítulo I** aborda la HCC, las bases electromagnéticas y las anomalías de las Pioneer y Flyby; el **Capítulo II** detalla la implementación y los resultados del Experimento I (medición directa de \\(c\\)); el **Capítulo III** expone exhaustivamente la arquitectura, construcción y resultados del Experimento II (circuito resonante modular RLC); el **Capítulo IV** plantea las propuestas de experimentación avanzada; y se finaliza con las **Conclusiones**, **Referencias** y **Apéndices** técnicos.


